
\documentclass{article}
\usepackage{colortbl}
\usepackage{makecell}
\usepackage{multirow}
\usepackage{supertabular}

\begin{document}

\newcounter{utterance}

\centering \large Interaction Transcript for game `ifeval', experiment `ifeval', episode 367 with refineonly{-}Qwen3.5{-}9B.
\vspace{24pt}

{ \footnotesize  \setcounter{utterance}{1}
\setlength{\tabcolsep}{0pt}
\begin{supertabular}{c@{$\;$}|p{.15\linewidth}@{}p{.15\linewidth}p{.15\linewidth}p{.15\linewidth}p{.15\linewidth}p{.15\linewidth}}
    \# & \multicolumn{2}{c}{Player} && \multicolumn{2}{c}{Game Master} \\
    \hline

    \theutterance \stepcounter{utterance}  
    & & & \multicolumn{4}{p{0.6\linewidth}}{
        \cellcolor[rgb]{0.9,0.9,0.9}{
            \makecell[{{p{\linewidth}}}]{
                \texttt{\tiny{[P1$\langle$GM]}}
                \texttt{Expand the riddle into a story with a funny tone:} \\
\\ 
\texttt{What can you catch but not throw?} \\
\texttt{A cold} \\
\\ 
\texttt{Use * to highlight at least 2 sections in your text. For example: *this is a highlighted text section*.} \\
            }
        }
    }
     \\ \\

    \theutterance \stepcounter{utterance}  
    & & \multicolumn{4}{p{0.6\linewidth}}{
        \cellcolor[rgb]{0.95,0.95,0.95}{
            \makecell[{{p{\linewidth}}}]{
                \texttt{\tiny{[P1$\leftrightarrow$P]}}
                \texttt{{[}DRAFT{]}} \\
\texttt{Barnaby was a man who took the phrase "catch a cold" quite literally. While everyone else treated the common cold as a metaphorical sneeze that happens to you, Barnaby believed it was a literal projectile waiting to be intercepted.} \\
\\ 
\texttt{One particularly chilly Tuesday, he stood in the middle of the grocery store aisle, wearing a tiny, ridiculous net made of fishing line and a pair of oversized goggles. He was scanning the air with intense focus, muttering to himself, "Just need to snag that one before it hits my nose."} \\
\\ 
\texttt{Suddenly, a shivering customer walked past him, sneezing violently. *Barnaby lunged forward, his net flailing wildly, trying to scoop up the invisible virus floating in the air.* "Gotcha!" he cheered, though he only managed to catch a stray piece of lint from his own sweater.} \\
\\ 
\texttt{His neighbor, Mrs. Gable, rolled her eyes so hard it looked like she was trying to see the back of her own head. "Barnaby," she sighed, adjusting her scarf, "you can't throw a cold, and you certainly can't catch it like a baseball. It's just a bug that decides to live inside you."} \\
\\ 
\texttt{"But the riddle says 'What can you catch but not throw?'!" Barnaby insisted, pulling the net down over his head like a helmet. "It has to be a cold! I'm going to go home and practice my throwing technique so I can finally get rid of it."} \\
\\ 
\texttt{*He spent the next three days lying on his couch, convinced he was training to launch the cold out of his body like a cannonball, while actually just getting progressively sicker.* By the time he finally admitted defeat and took some medicine, he realized the joke was on him: he had successfully "caught" the cold, but since he couldn't throw it away, he was stuck with it for a whole week.} \\
            }
        }
    }
    & \\ \\

    \theutterance \stepcounter{utterance}  
    & \multicolumn{4}{p{0.6\linewidth}}{
        \cellcolor[rgb]{0.561,0.737,0.561}{
            \makecell[{{p{\linewidth}}}]{
                \texttt{\tiny{[P1$\rangle$GM]}}
                \texttt{Barnaby was a man who took the phrase "catch a cold" quite literally. While everyone else treated the common cold as a metaphorical sneeze that happens to you, Barnaby believed it was a literal projectile waiting to be intercepted.} \\
\\ 
\texttt{One particularly chilly Tuesday, he stood in the middle of the grocery store aisle, wearing a tiny, ridiculous net made of fishing line and a pair of oversized goggles. He was scanning the air with intense focus, muttering to himself, "Just need to snag that one before it hits my nose."} \\
\\ 
\texttt{Suddenly, a shivering customer walked past him, sneezing violently. *Barnaby lunged forward, his net flailing wildly, trying to scoop up the invisible virus floating in the air.* "Gotcha!" he cheered, though he only managed to catch a stray piece of lint from his own sweater.} \\
\\ 
\texttt{His neighbor, Mrs. Gable, rolled her eyes so hard it looked like she was trying to see the back of her own head. "Barnaby," she sighed, adjusting her scarf, "you can't throw a cold, and you certainly can't catch it like a baseball. It's just a bug that decides to live inside you."} \\
\\ 
\texttt{"But the riddle says 'What can you catch but not throw?'!" Barnaby insisted, pulling the net down over his head like a helmet. "It has to be a cold! I'm going to go home and practice my throwing technique so I can finally get rid of it."} \\
\\ 
\texttt{*He spent the next three days lying on his couch, convinced he was training to launch the cold out of his body like a cannonball, while actually just getting progressively sicker.* By the time he finally admitted defeat and took some medicine, he realized the joke was on him: he had successfully "caught" the cold, but since he couldn't throw it away, he was stuck with it for a whole week.} \\
            }
        }
    }
    & & \\ \\

    \theutterance \stepcounter{utterance}  
    & & & \multicolumn{2}{p{0.3\linewidth}}{
        \cellcolor[rgb]{0.9,0.9,0.9}{
            \makecell[{{p{\linewidth}}}]{
                \texttt{\tiny{[GM$|$GM]}}
                {'detectable_format:number_highlighted_sections': {'num_highlights': 2}}
            }
        }
    }
    & & \\ \\

    \theutterance \stepcounter{utterance}  
    & & & \multicolumn{2}{p{0.3\linewidth}}{
        \cellcolor[rgb]{0.9,0.9,0.9}{
            \makecell[{{p{\linewidth}}}]{
                \texttt{\tiny{[GM$|$GM]}}
                \texttt{game\_result = WIN} \\
            }
        }
    }
    & & \\ \\

\end{supertabular}
}

\end{document}
